\clearpage
\section{Verification and Testing}

The implementation is checked through both \textbf{unit tests} and cluster-level
integration tests.

\subsection{Automated tests}

\begin{itemize}
    \item \texttt{ControlUnitActorTest} checks the state machine, full and partial arming, timers, siren control, and recovery rules.
    \item \texttt{SensorActorTest} and \texttt{KeypadActorTest} check message forwarding and behavior when the control unit is temporarily unavailable.
    \item \texttt{ClusteredSystemTest} verifies three-node cluster formation, remote message exchange, and recreation of the control unit in recovery mode.
    \item \texttt{NodeStartupTest} validates CLI parsing and seed-node configuration generation.
\end{itemize}

These tests are useful not only because they confirm correctness, but because
they document the expected \textbf{distributed behavior} of the final system.

\subsection{What is actually verified}

The automated suite verifies the properties that matter most for this
assignment:

\begin{itemize}
    \item \textbf{preservation} of the state transitions already tested in assignment 3;
    \item \textbf{three-node clustered startup} with distinct control-unit, keypad, and sensor roles;
    \item \textbf{message-only communication} between distributed actors;
    \item \textbf{safe recovery} after control-unit recreation;
    \item \textbf{serialization} of remote messages and snapshots.
\end{itemize}

In other words, the tests are centered on \textbf{behavioral correctness}, not on
advanced cluster features such as persistence, replication, or throughput.
